\section{Related work}
\label{sec:related}

\textbf{Fast-weight memory.} Writing transient associations into
multiplicative weights dates to \citet{schmidhuber1992learning}, revived
for attention-era models by \citet{ba2016using}. \citet{schlag2021linear}
identified linear attention \citep{katharopoulos2020transformers} as a
fast-weight programmer and introduced the delta-rule update our contender
uses; \citet{yang2024parallelizing} made that rule trainable at scale.
Selective state-space models \citep{gu2023mamba} pursue the same
constant-state property with a different recurrence, gated delta rules
refine the contender's own update \citep{yang2024gated}, and test-time
memorization modules \citep{behrouz2024titans} share this paper's
motivating contrast between a fixed-size memory and a growing KV cache.
None of these works
reports a parameter- and token-matched comparison against both a
transformer and a structure-removed vector control on episode-restricted
recall, with a causal state intervention and a byte-pinned cache-cap
protocol; that combination is this paper's contribution.

\textbf{Recall as the diagnostic.} \citet{arora2023zoology} established
multi-query associative recall as the axis separating attention from
efficient alternatives, with the opposite headline: in their regime,
attention wins and recurrent models lag. Our result is not a
contradiction but a different operating point: a small matched budget
(14M-class models, 20{,}000 steps from scratch), higher load per state
dimension ($K/d = 0.5$), an episode-restricted argmax metric, and a task
distribution built for exact matched-budget control. We also report where
our separation dies ($K/d = 0.75$, Section~\ref{sec:scope}), which their
framing predicts it must somewhere.

\textbf{KV-cache reduction.} Sink-plus-window eviction
\citep{xiao2023efficient} and heavy-hitter selection \citep{zhang2023h2o}
compress the cache of a model that already performs the task, asking how
few bytes preserve its behavior. We use the sink-plus-FIFO mechanism of
\citet{xiao2023efficient} in the opposite role, as an evaluation
instrument that pins the baseline's cache to a byte budget matched against
a fixed-size state, and our question is prior to theirs: whether the
matched attention model learns the behavior at all.

\textbf{Associative-memory capacity.} \citet{nichani2024understanding}
analyze factual recall in transformers via associative memories and prove
that under argmax decoding a rank-one outer-product store supports on the
order of $d$ associations. We import their result as a standing caveat on
what accuracy metrics can claim (Section~\ref{sec:setup}); our
contribution is empirical and comparative, not a capacity bound.
